\chapter{Affine varieties}

\lecture{1}{Thu 03 Sep 2026 11:07}{Syllabus day}

We will follow Andreas Gathmann's 2002 notes \cite{Gathmann2002}, and will suppliment them with the various books on Algebraic Geometry, such as \cite{vakil} and Hartshorne.

Algebraic geometry is essentially the study of the geometry of solution sets to systems of polynomial equations, using the tools of commutative algebra. We set up a dictionary to translate between geometric and algebraic properties, which holds more generally than just polynomial rings. Given a sufficiently nice space $X$, we obtain a ring $C(X)$ of $\mathbb{C} $-valued continuous functions. There's a number of things we can do here. For example, given a point $x\in X$, there is an evaluation map $\varepsilon _x : C(X) \to \mathbb{C} $ which ends up being a ring homomorphism. Since the space is nice, we can take $\varepsilon _x$ to be surjective. Thus, $\Ker \varepsilon _x$ is a maximal ideal of $C(X)$ due to the first isomorphism theorem $C(X) / \Ker \varepsilon _x \cong \mathbb{C} $. Thus, something as simple as a point $x\in X$ manifests as a maximal ideal in the ring $C(X)$.

More generally, let $Z\subseteq X$ be closed. I can define an ideal $I_Z \subseteq C(X)$ of all functions $f\in C(X)$ that vanish on $Z$. Thus, arbitrary closed sets give us an ideal that is not necessarily maximal.

Consider the space $X$ given by the union of the $x$- and $y$-axes inside of $\mathbb{C} ^2$. Letting $f=y$ and $g=x$, we see that $f$ vanishes on the $y$-axis, and $g$ vanishes on the $x$-axis. Then $fg=xy$ is identically zero on $X$. Thus, the ring $C(X)$ has zero divisors.

\section{Algebraic sets and the Zariski topology}

\begin{definition}\label{dfn:1.1.1}
	Let $k$ be a field. Define the \emph{$n$-dimensional affine space over $k$}, denoted $\mathbb{A} _n^k$, as
	\[
		\mathbb{A}^n_k \coloneqq \left\{ (x_1, \ldots ,x_n)\mid x_i\in k \forall i \right\} 
	.\]
\end{definition}
The intuition for affine space is that we try to think of it as an $n$-dimensional vector space without an origin. We will also want to consider non-linear systems of equations over affine space, which will lead us to topologize it in a strange way. This topology is what differentiates $\mathbb{A} _n^k$ from $k^n$ in the case where $k$ is a topological field.

For now, note that we can view $k[x_1, \ldots ,x_n]$ as $k$-valued polynomial functions on $\mathbb{A} _k^n$. Given a subset $S\subseteq k[x_1, \ldots ,x_n]$, we want to find the \emph{solution set} or \emph{vanishing locus}
\[
	Z(S) = \left\{ x\in \mathbb{A} _k^n \mid \forall p\in S, p(x)=0 \right\} 
.\]

We can make some observabtions. If $S_1 \subseteq S_2 \subseteq k[x_1, \ldots ,x_n]$, then since $S_2$ applies more relations, $Z(S_2)\subseteq Z(S_1)$. Thus $Z$ is inclusion reversing. The next observation is that it would be more natural to consider ideals. This is further motivated by the fact that if $(S)$ is the ideal generated by $S$, then we can show $Z(S) = Z((S))$. To prove this, first note that since $S\subseteq (S)$, we have $Z((S))\subseteq Z(S)$. For the converse, let $x\in Z(S)$. Note that an element of $(S)$ is of the form
\[
	\sum_{i} g_ip_i,\qquad g_i\in k[x_1, \ldots ,x_n],\qquad p_i\in S
.\]
Since $p_i(x)=0$ for all $i$, this sum must vanish at $x$. If $S=\left\{ f_1, \ldots ,f_r \right\} $ is finite, then we often write $Z(f_1, \ldots ,f_r)$. This discussion leads us to give sets of the form $Z(S)$ a special name.

\begin{definition}\label{dfn:1.1.2}
	An \emph{algebraic set} is a subset of $ \mathbb{A} ^n_k$ of the form $Z(S)$, where $S\subseteq k[x_1, \ldots ,x_n]$.
\end{definition}

For example, $\mathbb{A} ^n_k = Z(0)$, so affine space is an algebraic set. The empty set is the vanishing locus of any nonzero constant function, so it is also an algebraic set. A sinle point $x\in\mathbb{A} _k^n$ is the vanishing locus of $Z(x_1 - a_1, \ldots ,x_n - a_n)$, so singletons are algebraic sets. More generally, any linear or affine linear subspace of $\mathbb{A} ^n_k$ is an algebraic set. For example, $x_2 - x_1$ cuts out a diagonal in the $x_1x_2$-plane. The circle is the vanishing locus of $(x_1^2 + x_2^2 - 1)$.


We will now define the topology we alluded to earlier, referred to as the \emph{Zariski topology}. The idea is that we want to topologize $\mathbb{A} _n^k$ such that the closed sets are precisely the algebraic sets.

\begin{proposition}\label{prp:1.1.3}
	There is a topology on $\mathbb{A} _k^n$ whose closed sets are precisely the algebraic sets.
\end{proposition}
\begin{proof}
	We must show the following.
	\begin{enumerate}
		\item $\varnothing ,\mathbb{A}^n_k$ are closed,
		\item a finite union of closed sets is closed, and
		\item arbitrary intersections of closed sets are closed.
	\end{enumerate}
	We'be already shown property 1. The remaining results follow immediately from the following lemma -- since the relevant constructions end up being vanishing loci of ideals, the result follows.
\end{proof}

\begin{lemma}\label{lma:1.1.4}
	Let $\left\{ I_j \right\} _{j\in J} $ be ideals of $k[x_1, \ldots ,x_n]$. Then
	\[
		\bigcap_{j\in J} Z(I_j) = Z\left( \sum_{j\in J} I_j \right) 
	.\]
	Additionally, if $J=[n]$ is finite, then
	\[
		\bigcup_{j=0}^n Z(I_j) = Z\left( \prod_{j=0}^n I_j \right) = Z\left( \bigcap_{j=0}^n I_j \right) 
	.\]
\end{lemma}
\begin{proof}
	Recall a sum of ideals is all finite linear combinations in those ideals, and the product of ideals is given by all monomials using the normal ring product (so thing $gf_1\cdots f_n$).
\end{proof}

This topology is different from the usual topology on $\mathbb{A} _k^n$ (when $k=\mathbb{R} ,\mathbb{C} $). We call this the \emph{Zariski topology}. Whenever we refer to affine space, we implicitly take the given topology to be the Zariski topology.

Let's look at some closed sets of $\mathbb{A} _k^1$. Let $f$ be a polynomial over an algebraically closed field $k$. Then $f$ splits as $(x-a_1)\cdots (x-a_r)$, so $Z(f)=\left\{ a_1, \ldots ,a_r \right\} $. Thus since $k$ is algebraically closed, the closed sets of 1 dimensional affine space are finite collections of points. Thus the opens are FUCKING huge in the Zariski topology, and the Zariski topology is \emph{far} courser than the analytic topology. In fact, any two nonempty opens will intersect in $\mathbb{A} _k^1$ (and I think in any affine space). Of course, closed sets in the Zariski topology can have uncountably many points in higher dimensions (recall all our earlier examples).

Finally, we recall Hilbert's basis theorem, which will be useful later on in the course.

\begin{theorem}[Hilbert's basis theorem]\label{thm:hilbertbasis}
	Let $R$ be a Noetherian ring. Then $R[x]$ is a Noetherian ring.
\end{theorem}


